\section*{Chapitre \ref{ch:b3:sequence-stratigraphy} --- Stratigraphie séquentielle}

\begin{solution}{exo:b3:sequence-stratigraphy:1}
$A = 0.5 + 0.1 = \qty{0.6}{m/ka}$~; en \qty{20}{ka}, \qty{12}{m} d'espace
d'accommodation. Avec la mer qui baisse, $A = -0.3 + 0.1 = \qty{-0.2}{m/ka}$~:
une baisse relative, ralentie par la subsidence.
\end{solution}

\begin{solution}{exo:b3:sequence-stratigraphy:2}
(a) \omterm{def:b3:sequence-stratigraphy:tract}{Cortège de chute} (\omterm{def:b3:sequence-stratigraphy:regression}{régression forcée}). (b) \omterm{def:b3:sequence-stratigraphy:tract}{Cortège transgressif}. (c)
\omterm{def:b3:sequence-stratigraphy:tract}{Cortège de haut niveau}. (d) \omterm{def:b3:sequence-stratigraphy:tract}{Cortège de bas niveau} (\omterm{def:b3:sequence-stratigraphy:regression}{régression normale} après la baisse).
\end{solution}

\begin{solution}{exo:b3:sequence-stratigraphy:3}
Chaque \omterm{def:b3:sequence-stratigraphy:parasequence}{paraséquence} se termine moins profondément que celle qui la précède~: un
empilement progradant, le rivage se déplaçant vers le large, comme dans un
\omterm{def:b3:sequence-stratigraphy:tract}{cortège de haut niveau} ou de bas niveau. L'ordre inverse, chacune se terminant
plus fine et plus profonde, est rétrogradant~: c'est un \omterm{def:b3:sequence-stratigraphy:tract}{cortège transgressif}.
\end{solution}

\begin{solution}{exo:b3:sequence-stratigraphy:4}
Cinq couplets par faisceau correspondent au rapport entre l'excentricité courte
(environ \qty{100}{ka}) et la précession (environ \qty{20}{ka})~: chaque
couplet est un cycle de précession, chaque faisceau un cycle d'excentricité. Le
taux vaut $1.6/100 = \qty{1.6}{cm/ka}$, et un couplet de \qty{0.32}{m} dure
environ \qty{20}{ka}, ce qui est cohérent.
\end{solution}

\begin{solution}{exo:b3:sequence-stratigraphy:5}
$\mathrm dx_s/\mathrm dt = (4000 - 0.15\times20\,000)/50 = \qty{20}{m/ka}$,
vers le large. Le rivage reste immobile lorsque $AW = Q$~: $A = 4000/20\,000 =
\qty{0.2}{m/ka}$.
\end{solution}

\begin{solution}{exo:b3:sequence-stratigraphy:6}
$\mathrm dx_s/\mathrm dt = 4000/50 + 0.2/0.01 = 80 + 20 = \qty{100}{m/ka}$~;
la seule baisse explique \qty{20}{m/ka}, soit un cinquième.
\end{solution}

\begin{solution}{exo:b3:sequence-stratigraphy:7}
Il recule lorsque $Q < AW = 0.4\times25\,000 = \qty{10000}{m^3/ka}$ par mètre.
Avec $Q = 6000$~: $(6000 - 10\,000)/10 = \qty{-400}{m/ka}$, soit un recul de
\qty{0.4}{km} tous les mille ans~: un avant-côte bas produit une transgression
rapide.
\end{solution}

\begin{solution}{exo:b3:sequence-stratigraphy:8}
$1/94.9 - 1/123.9 = \qty{0.00247}{cycles/ka}$~: il faut un enregistrement plus
long que $1/0.00247 = \qty{406}{ka}$ pour les séparer, et \qty{600}{ka}
(résolution 0.00167) suffisent. $1/94.9 - 1/98.9 = 0.00042$~: il faut environ
\qty{2400}{ka}, soit quatre fois la durée de l'enregistrement.
\end{solution}

\begin{solution}{exo:b3:sequence-stratigraphy:9}
$16.4/(20\times40.98) = \qty{0.0200}{m/ka}$, \qty{2.0}{cm/ka}~; les lits de
cendres sont distants de $3.1/0.020 = \qty{155}{ka}$.
\end{solution}

\begin{solution}{exo:b3:sequence-stratigraphy:10}
(a) $\tan\alpha = -0.4/200 = -0.002$~: descente à \ang{0.11}. (b)
$0.2/50 = 0.004$~: remontée à \ang{0.23}. Une section sismique couvre des
dizaines de kilomètres avec une exagération verticale de dizaines à centaines,
si bien que l'angle se voit~; une falaise de quelques centaines de mètres de
long montre une remontée d'un mètre ou moins, noyée dans la variation des lits.
\end{solution}

\begin{solution}{exo:b3:sequence-stratigraphy:11}
$1\,296\,000/54.725282 = \qty{23.68}{ka}$ ($g_5$)~; $/57.924383 =
\qty{22.37}{ka}$ ($g_2$)~; $/68.377912 = \qty{18.95}{ka}$ ($g_4$)~;
$/67.834313 = \qty{19.11}{ka}$ ($g_3$). $g_2 - g_5 = 3.199101''$ par an,
\qty{405.1}{ka}~; $g_4 - g_5 = 13.65263''$, \qty{94.9}{ka}.
\end{solution}

\begin{solution}{exo:b3:sequence-stratigraphy:12}
Les deux donnent $0.3\times100 = \qty{30}{m}$ d'espace d'accommodation au même
rythme~; avec le même apport, la géométrie est identique, puisque les sédiments
ne répondent qu'au niveau marin relatif. Une variation eustatique est la même
partout~: la même séquence, de même âge, sur d'autres marges, désigne la
première histoire. Une histoire de subsidence reconstruite à partir de
l'histoire thermique et de charge du bassin (backstripping), ou un
enregistrement eustatique indépendant tel que les isotopes de l'oxygène, les sépare également.
\end{solution}

\begin{solution}{pb:b3:sequence-stratigraphy:1}
\textbf{1.} $\dot E = 30\times2\pi/400\times\sin(2\pi t/400)$~: au maximum
\qty{0.471}{m/ka}, croissant à \qty{100}{ka} et \qty{500}{ka}, décroissant à
\qty{300}{ka}.

\textbf{2.} \qty{0.020}{m/ka} à la charnière, 0.037 à \qty{20}{km}, 0.062 à
\qty{50}{km}, \qty{0.120}{m/ka} à \qty{120}{km}.

\textbf{3.} $A < 0$ lorsque $\sin(2\pi t/400) < -0.0346/0.471 = -0.0734$~: $t =
200 + (400/2\pi)\arcsin0.0734 = \qty{204.7}{ka}$.

\textbf{4.} Avec $\sigma = 0.0617$~: $t = 400 - 63.66\arcsin0.131 =
\qty{391.6}{ka}$.

\textbf{5.} Environ \qty{187}{ka} contre 200~: la baisse eustatique doit
d'abord dépasser la subsidence, et cesse de la vaincre avant le bas niveau.

\textbf{6.} $0.045\times400 = \qty{18}{m}$, contre une amplitude eustatique de
\qty{60}{m}~: un tiers.

\textbf{7.} Vers l'intérieur des terres de 50 à environ \qty{120}{ka}, vers le large jusqu'à environ
\qty{440}{ka}, vers l'intérieur jusqu'à environ \qty{560}{ka}, de nouveau vers le large jusqu'à
\qty{600}{ka} (les \omterm{def:b3:seismology-imaging:ray-parameter}{points de retournement} du modèle sont 118.5, 439.5 et
\qty{559}{ka}~; un tableau tous les \qty{50}{ka} ne fait que les encadrer).

\textbf{8.} $(6000 - 0.30\times15\,000)/40 = \qty{37.5}{m/ka}$, vers le large.

\textbf{9.} $6000/15\,000 = \qty{0.40}{m/ka}$.

\textbf{10.} $6000/35.7 + 0.424/0.012 = 168.1 + 35.3 = \qty{203}{m/ka}$,
à moins de \qty{0.5}{\percent} du modèle.

\textbf{11.} L'apport~: \qty{83}{\percent} de la vitesse.

\textbf{12.} $(-42.7 - 16.3)/(50.7 - 17.4) = \qty{-1.77}{m/km}$, \ang{0.10}~O.
Depuis son instant, le premier point s'est affaissé de $0.0345\times395.5 =
\qty{13.6}{m}$, le second de $0.0623\times208.5 = \qty{13.0}{m}$~: presque
la même valeur, de sorte que, aux altitudes actuelles, la descente, \qty{59.0}{m}, ne diffère
de la baisse eustatique que de \qty{0.6}{m}. La baisse relative du rivage à
ce moment, soit l'eustatisme moins la subsidence pendant la baisse, était plus faible.

\textbf{13.} 150 et 200~: haut niveau~; 250, 300, 350~: \omterm{def:b3:sequence-stratigraphy:tract}{cortège de chute}~; 400~:
bas niveau~; 450, 500, 550~: transgressif (le rivage recule juste après
440)~; 600~: haut niveau.

\textbf{14.} \qty{204.7}{ka} et \qty{391.6}{ka} contre 204.5 et 391.5~: les mêmes,
à la précision du pas du modèle, \qty{0.5}{ka}.

\textbf{15.} Sur la plate-forme, la surface exposée pendant la baisse~: le sommet des
dépôts de haut niveau, de la charnière jusqu'à \qty{17}{km}, se prolonge au-dessus des
foresets de \omterm{def:b3:sequence-stratigraphy:tract}{cortège de chute} en escalier~; dans le bassin, la concordance corrélative
à la base ou au sommet du \omterm{def:b3:sequence-stratigraphy:tract}{cortège de chute}, par convention. Une vraie plate-forme
serait entaillée par des rivières en \omterm{def:b3:sequence-stratigraphy:condensed}{vallées incisées}, altérée en sols et
en partie érodée.

\textbf{16.} $(18.5 - 52.8)/(559 - 439.5) = \qty{-0.29}{km/ka}$~; la
\omterm{def:b3:sequence-stratigraphy:surfaces}{surface d'inondation maximale} se forme à sa fin.

\textbf{17.} Le rivage balaie rapidement une plaine côtière presque plate~:
l'apport est piégé près de lui, réparti sur une large zone, et chaque endroit
en reçoit peu. Au-delà du pied de pente, au-delà d'environ \qty{55}{km}, le fond
marin est affamé~: une \omterm{def:b3:sequence-stratigraphy:condensed}{section condensée} s'y forme.

\textbf{18.} Résolution $1/30 = \qty{0.033}{cycles/m}$~; Nyquist $1/(2\times
0.01) = \qty{50}{cycles/m}$.

\textbf{19.} $7.5/2.0 = 3.75$, 1 et $0.81/2.0 = 0.405$~; les rapports orbitaux
sont $405/100 = 4.05$ et $41/100 = 0.41$.

\textbf{20.} \qty{2.0}{m} pour \qty{100}{ka}~: \qty{2.0}{cm/ka}~; $0.81/41 =
\qty{1.98}{cm/ka}$~; $7.5/405 = \qty{1.85}{cm/ka}$. Les longueurs d'onde que le
\omterm{def:b3:sequence-stratigraphy:periodogram}{périodogramme} peut montrer près de \qty{8}{m} sont 10, 7.5 et \qty{6}{m}~: la
vraie valeur \qty{8.1}{m} tombe entre deux d'entre elles.

\textbf{21.} $30/0.020 = \qty{1500}{ka}$.

\textbf{22.} $23.7\times0.020 = \qty{0.474}{m}$. Échantillonnée tous les
\qty{25}{cm}, la fréquence de Nyquist vaut 2 cycles par mètre, soit une longueur d'onde de
\qty{0.5}{m}~: la précession, à 2.11 cycles par mètre, lui est supérieure et
se replie en $4 - 2.11 = 1.89$, une fausse longueur d'onde de \qty{0.53}{m}.

\textbf{23.} Un pic isolé s'ajuste à une période quelconque pour un taux quelconque,
quelle que soit son origine~; trois pics dont les rapports correspondent aux rapports
orbitaux sont indépendants du taux, qui doit alors ressortir identique pour chacun.

\textbf{24.} Le rivage passe de \qty{17.37}{km} au début de la baisse à
\qty{52.84}{km} à la fin de la régression de bas niveau~: un déplacement
maximal vers le large de \textbf{\qty[text-series-to-math=true]{35.5}{km}}.
\end{solution}
